\subsection{Configuration and braid conventions}

The ordered configuration space of $n$ points in a surface $S$ is
\[
 \Conf_n(S)=\{(z_1,\ldots,z_n)\in S^n:z_i\ne z_j\text{ for }i\ne j\},
\]
and $\UConf_n(S)=\Conf_n(S)/S_n$.  For the plane,
$\pi_1(\UConf_n(\C))=B_n$.  Our generators $\sigma_i$ exchange the $i$th and
$(i+1)$st points by a positive half-twist and satisfy
\begin{align*}
 \sigma_i\sigma_j&=\sigma_j\sigma_i &&(|i-j|\ge2),\\
 \sigma_i\sigma_{i+1}\sigma_i
 &=\sigma_{i+1}\sigma_i\sigma_{i+1}.&&
\end{align*}
The quotient $B_n\to S_n$ remembers only endpoint permutation, and the
abelianization $B_n\to\Z$ sends every $\sigma_i$ to $1$.

The identification \eqref{eq:vieta-configuration} sends an unordered set of roots
to the elementary symmetric coefficients of its monic polynomial.  It is a
homeomorphism, with smooth complex-manifold structure on the square-free locus.
The classical configuration-space fibrations supply one standard route to the
fundamental-group computation \cite{FadellNeuwirth1962ConfigurationSpaces}.

\subsection{Properness of a finite root covering}

Let $p:Y\to B$ be a finite-sheeted covering with $B$ locally compact Hausdorff.
For a compact set $K\subset B$, cover $K$ by finitely many relatively compact
evenly covered neighborhoods.  The closure of each sufficiently small sheet is
homeomorphic to a compact closure downstairs.  A finite union contains
$p^{-1}(K)$ as a closed subset, so $p^{-1}(K)$ is compact.  Thus $p$ is proper.
Applied to the root incidence, this proof uses square-freeness locally and the
fixed finite degree globally.  Monicity prevents a root from disappearing to
infinity while the coefficient remains at a fixed point of $B$.

\subsection{Gauge calculation for logarithmic roots}

Let $z_i(s)$ be the continuation path beginning at the $i$th labeled root.  If
$z_i(1)=z_{\sigma(i)}(0)$ and $\omega_i$ is its logarithmic lift, then
\[
 \omega_i(1)-\omega_{\sigma(i)}(0)\in2\pi i\Z.
\]
This defines $k_i$.  Shift initial lifts by $m_j\in\Z$.  Unique path lifting gives
\[
 \omega_i'(s)=\omega_i(s)+2\pi im_i,
\]
and hence
\begin{align*}
 \omega_i'(1)
 &=\omega_{\sigma(i)}(0)+2\pi i(k_i+m_i)\\
 &=\omega'_{\sigma(i)}(0)
   +2\pi i(k_i+m_i-m_{\sigma(i)}).
\end{align*}
This proves \eqref{eq:log-gauge-law}.  If
$i_1\mapsto i_2\mapsto\cdots\mapsto i_\ell\mapsto i_1$ is a permutation cycle,
then
\[
 \sum_{j=1}^\ell
 (m_{i_j}-m_{\sigma(i_j)})=0.
\]
The cycle sum of deck shifts is therefore lift-gauge invariant, whereas each
summand is not.  Relabeling the roots permutes the cycles, so the unordered
multiset of cycle sums is canonical.

The total sum has an algebraic expression.  If the roots are $z_i(b)$, then
\[
 \prod_i z_i(b)=(-1)^n c_n(b).
\]
The winding of the constant coefficient around zero equals the sum of the root
windings and hence $\sum_i k_i$.

\subsection{Helical trace on the parameter-position torus}

On $S^1_\theta\times S^1_y$, the equation
$ny-m\theta\in\pi\Z$ is the inverse image of the two values $0$ and $\pi$ under
the group homomorphism
\[
 (\theta,y)\longmapsto ny-m\theta\pmod{2\pi}.
\]
The kernel of the homomorphism with coefficient vector $(-m,n)$ has
$d=\gcd(n,m)$ components.  Each of the two inverse images is a translate of that
kernel, giving $2d$ components in total.  A primitive parameterization of any
component is
\[
 s\longmapsto
 \left(\theta_0+\frac nds,\,
       y_0+\frac mds\right),
 \qquad s\in\R/(2\pi\Z),
\]
which proves \eqref{eq:helical-boundary-class}.  With the ordered torus-link
convention declared in Section~\ref{sec:regular-tubes}, the union is $T(2n,2m)$.

\subsection{The sign of the simple half-twist}

For $w^2=\tau$ on the positively oriented boundary
$\tau=\varepsilon e^{i\theta}$, choose at $\theta=0$ the ordering
$(-\sqrt\varepsilon,+\sqrt\varepsilon)$.  As $\theta$ increases, both points turn
counterclockwise through angle $\pi$ and exchange.  Under the convention that a
counterclockwise exchange in the root plane is positive, the braid is
$\sigma_1$.  Reversing the boundary orientation gives $\sigma_1^{-1}$.  Changing
the local coordinate $w$ by an orientation-preserving complex affine map does not
change this sign.

This local sign does not determine a global filling.  If a parameter disc meets the
discriminant in several transverse points, the boundary braid is a product of
conjugates of signed half-twists, with order determined by a choice of paths from a
basepoint.  Holomorphic dependence restricts the possible signs and factorizations;
the quasipositive boundary theorem is the relevant classical constraint
\cite{Rudolph1983AlgebraicFunctions}.

\subsection{Closure equivalence}

Alexander's theorem says that every oriented link in $S^3$ has a closed-braid
representative.  Markov's theorem says that two braids have isotopic oriented
closures precisely after a sequence generated by braid conjugation and positive or
negative stabilization, allowing changes of strand number.  Accordingly, a
quantity on $B_n$ is not a classical link invariant until the three identities in
\eqref{eq:markov-gate}, with its chosen normalization, are proved.  If the braid
axis is retained, one obtains an annular invariant under a smaller equivalence
relation; it must not be relabeled an ordinary $S^3$ invariant.
